\section{Background and Related Work}\label{sec:background}

\subsection{Man-in-the-browser and real-time phishing}
A man-in-the-browser Trojan sits inside the victim's browser and intercepts the
calls between the user interface and the network, so it sees the transaction the
user intends and can substitute a different one before it reaches the bank,
while keeping the displayed page unchanged. A real-time phishing proxy achieves
the same position from the network side: the victim interacts with an
attacker-controlled copy of the bank site, and the attacker relays and rewrites
the traffic. Because both end up between the user's intent and the bank's record
of the transaction, we treat them as one capability in our threat model
(\cref{sec:threat}).

\subsection{Transaction binding and the trusted display}
The idea of binding the authorisation to the specific transaction, and of
showing the transaction on a trusted device so the user signs what they see
(WYSIWYS), is well established in both research and deployment. Hiltgen,
Kramp and Weigold surveyed secure internet-banking authentication and argued for
transaction-specific confirmation~\cite{hiltgen2006}; the Zurich Trusted
Information Channel (ZTIC) realised a trusted device that displays and confirms
transaction data~\cite{ztic2008}. In deployment, chipTAN and photoTAN schemes
derive a transaction code from the transaction data, and the challenge-response
standard OCRA includes a signature/transaction mode (RFC~6287, signature
mode)~\cite{rfc6287}. Regulation has made the principle mandatory for European
payment service providers: the ``dynamic linking'' requirement of Commission
Delegated Regulation (EU)~2018/389, Article~5, requires the authentication code
to be specific to the amount and the payee~\cite{eu2018389}. More recent web
standards, such as W3C Secure Payment Confirmation, carry transaction details
into the signed authenticator assertion~\cite{w3cspc}.\footnote{W3C Secure
Payment Confirmation is, at the time of writing, a Candidate Recommendation
Draft rather than a full Recommendation.}

Against this background, the four modes in this report are not new designs.
\texttt{baseline} is RFC~4226/6238 HOTP/TOTP~\cite{rfc4226,rfc6238};
\texttt{hmac} is an OCRA-\emph{style} code (it uses HMAC-SHA256 over a canonical
transaction with a domain-separation tag, but it is \emph{not} standards-exact
OCRA: there is no OCRASuite string and no RFC~6287 \texttt{DataInput} layout);
\texttt{sig}/\texttt{sig\_blind} apply an ECDSA~P-256 signature (FIPS~186-5) over
the same canonical transaction. The contribution is the controlled comparison
and the blind-signing experiment, not the mechanisms.

\subsection{Do users actually check?}
WYSIWYS only helps if the user reads the trusted display and notices a mismatch.
This assumption is known to be optimistic. AlZomai, AlFayyadh, J\o sang and
McCullagh ran a simulated-bank user study of transaction
authorisation~\cite{alzomai2008}, and Haupert and Gabert surveyed real banking
customers on where they look for the signed transaction
details~\cite{haupert2019}. We use these works to bound the interpretation of
our attentive-user results (\cref{sec:eval,sec:discussion}); we do not reproduce
their specific detection rates here, and quote them only once read from the
original publications.

\subsection{Comparison with existing solutions}
\Cref{tab:compare} positions the four implemented modes against representative
existing approaches on the properties that matter for MitB defence. It is a
qualitative map, not a claim of novelty.

\begin{table}[htbp]
\centering
\caption{Where the four modes sit among transaction-authentication approaches.}
\label{tab:compare}
\small
\begin{tabularx}{\linewidth}{@{}lXccc@{}}
\toprule
Approach & Binds code to transaction? & Trusted display & Asymmetric & Standards \\
\midrule
Plain OTP / TOTP (\texttt{baseline})   & No  & No  & No  & RFC 4226/6238 \\
chipTAN / photoTAN                     & Yes & Yes & No  & vendor \\
OCRA signature mode                    & Yes & Yes & No  & RFC 6287 \\
HMAC transaction code (\texttt{hmac})  & Yes & Yes & No  & OCRA-style \\
Signature (\texttt{sig})               & Yes & Yes & Yes & FIPS 186-5 \\
Signature, no display (\texttt{sig\_blind}) & Yes & No & Yes & FIPS 186-5 \\
WebAuthn / W3C SPC                     & Yes (SPC) & Yes & Yes & W3C \\
\bottomrule
\end{tabularx}
\end{table}
